\chapter{Bilangan sampai 10\,000}\label{ch:g3:numbers}

Seribu, dua ribu, tiga ribu\,\dots\ bilangan besar dibangun dari
bilangan kecil, dengan satu gagasan yang indah: \emph{tempat} sebuah
angka menyatakan nilainya. Bab ini mengajarkan cara membaca, menulis,
membandingkan dan mengurutkan bilangan sampai $10\,000$.

\section{Membaca dan menulis bilangan}

\begin{definition}[Nilai tempat]\label{def:g3:numbers:place}
Sebuah bilangan ditulis dengan angka $0, 1, 2, \dots, 9$. Dibaca dari
kanan, tempatnya adalah: satuan, puluhan, ratusan, ribuan. Pada
$3\,254$:
\[
3\,254 = 3 \text{ ribuan} + 2 \text{ ratusan} + 5 \text{ puluhan}
+ 4 \text{ satuan}.
\]
\end{definition}

\begin{omfigure}
\begin{tikzpicture}[scale=0.9]
  \draw (0,0) grid[xstep=2.3, ystep=0.85] (9.2,1.7);
  \node[font=\small] at (1.15,1.275) {ribuan};
  \node[font=\small] at (3.45,1.275) {ratusan};
  \node[font=\small] at (5.75,1.275) {puluhan};
  \node[font=\small] at (8.05,1.275) {satuan};
  \node[font=\large, omDef] at (1.15,0.425) {$3$};
  \node[font=\large, omDef] at (3.45,0.425) {$2$};
  \node[font=\large, omDef] at (5.75,0.425) {$5$};
  \node[font=\large, omDef] at (8.05,0.425) {$4$};
\end{tikzpicture}

{\small Tabel nilai tempat dari $3\,254$, dibaca ``tiga ribu dua
ratus lima puluh empat''.}
\end{omfigure}

\begin{example}[Nol yang menjaga tempat]\label{ex:g3:numbers:zero}
``Empat ribu tujuh'' ditulis $4\,007$: angka $4$ di ribuan, angka
$7$ di satuan, dan \emph{\omterm{def:g1:counting:zero}{nol}} di antaranya untuk menahan setiap
angka di tempatnya yang benar. Tanpa \omterm{def:g1:counting:zero}{nol} itu, $47$ menjadi bilangan
yang sangat berbeda!
\end{example}

\begin{example}[Menukar]\label{ex:g3:numbers:exchange}
Sepuluh satuan menjadi satu puluhan; sepuluh puluhan menjadi satu
ratusan; sepuluh ratusan menjadi satu ribuan. Jadi $10$ ratusan
$= 1\,000$, dan $32$ puluhan $= 320$: bilangan $320$ memuat $32$
puluhan (bukan hanya angka $2$!).
\end{example}

\section{Membandingkan dan mengurutkan}

\begin{method}[Membandingkan dua bilangan]\label{met:g3:numbers:compare}
\begin{enumerate}
  \item Bilanglah angkanya: makin banyak angka berarti makin besar
        bilangannya ($1\,002 > 998$).
  \item Banyak angkanya sama: bandingkan angka demi angka mulai dari
        \emph{kiri}. Perbedaan pertama yang menentukan:
        $5\,614 > 5\,589$ karena $6 > 5$ pada ratusan.
  \item Tulislah jawabannya dengan tanda $<$ (lebih kecil daripada)
        atau $>$ (lebih besar daripada): sisi yang terbuka selalu
        menghadap bilangan yang lebih besar.
\end{enumerate}
\end{method}

\begin{omfigure}
\begin{tikzpicture}[scale=1.1]
  \draw[-Stealth] (-0.3,0) -- (10.5,0);
  \foreach \i in {0,1,...,10} \draw (\i,-0.08) -- (\i,0.08);
  \node[below, font=\small] at (0,-0.1) {$0$};
  \node[below, font=\small] at (5,-0.1) {$500$};
  \node[below, font=\small] at (10,-0.1) {$1000$};
  \fill[omThm] (2.4,0) circle (2pt) node[above, font=\small] {$240$};
  \fill[omDef] (6.9,0) circle (2pt) node[above, font=\small] {$690$};
  \fill[omMeth] (9.05,0) circle (2pt) node[above, font=\small] {$905$};
\end{tikzpicture}

{\small Bilangan pada garis berskala ratusan: makin ke kanan makin
besar. $240 < 690 < 905$.}
\end{omfigure}

\begin{example}\label{ex:g3:numbers:order}
Urutkan dari yang terkecil ke yang terbesar: $780$, $8\,700$, $807$,
$87$. Mulailah dari banyaknya angka: $87$ (dua), lalu $780$ dan
$807$ (tiga), lalu $8\,700$ (empat). Antara $780$ dan $807$:
bandingkan angka ratusannya, $7 < 8$, jadi $780 < 807$. Urutan
akhirnya:
\[
87 < 780 < 807 < 8\,700 .
\]
\end{example}

\section{Genap dan ganjil}

\begin{definition}[Genap, ganjil]\label{def:g3:numbers:evenodd}
Sebuah bilangan disebut \emph{genap}\index{bilangan genap} bila ia
dapat dibagi menjadi dua bagian utuh yang sama besar --- angka
satuannya $0$, $2$, $4$, $6$ atau $8$. Selain itu ia
\emph{ganjil}\index{bilangan ganjil} (angka satuan $1$, $3$, $5$,
$7$ atau $9$).
\end{definition}

\begin{example}
$46$ \omterm{def:g3:numbers:evenodd}{genap} ($46 = 23 + 23$); $47$ \omterm{def:g3:numbers:evenodd}{ganjil} (dua bagian yang sama akan
memerlukan separuh benda). Hanya angka \emph{terakhir} yang
menentukan: $3\,578$ \omterm{def:g3:numbers:evenodd}{genap}, $2\,341$ \omterm{def:g3:numbers:evenodd}{ganjil}.
\end{example}

\begin{omfigure}
\begin{tikzpicture}[scale=0.45]
  \foreach \i in {0,...,5} \fill[omDef] ({\i},0) circle (0.28);
  \foreach \i in {0,...,5} \fill[omDef] ({\i},1) circle (0.28);
  \node[below, font=\small] at (2.5,-0.7) {$12$: genap, dua baris sama};
  \begin{scope}[xshift=10.5cm]
  \foreach \i in {0,...,4} \fill[omThm] ({\i},0) circle (0.28);
  \foreach \i in {0,...,4} \fill[omThm] ({\i},1) circle (0.28);
  \fill[omThm] (5,0) circle (0.28);
  \draw[omExo, thick] (5,0) circle (0.5);
  \node[below, font=\small] at (2.5,-0.7) {$11$: ganjil, satu titik sendirian};
  \end{scope}
\end{tikzpicture}

{\small \omterm{def:g3:numbers:evenodd}{Bilangan genap} terbagi menjadi dua baris yang sama;
\omterm{def:g3:numbers:evenodd}{bilangan ganjil} selalu menyisakan satu yang sendirian.}
\end{omfigure}

\section{Latihan}

\begin{exercise}[$\star$]\label{exo:g3:numbers:1}
Tulislah dengan angka: ``dua ribu tiga ratus empat puluh lima'';
``enam ribu lima puluh''; ``sembilan ribu satu''.
\end{exercise}

\begin{exercise}[$\star$]\label{exo:g3:numbers:2}
Tulislah dengan kata: $1\,208$; \quad $4\,070$; \quad $9\,999$.
\end{exercise}

\begin{exercise}[$\star$]\label{exo:g3:numbers:3}
Pada $7\,382$: angka mana yang menempati tempat puluhan? Tempat
ribuan? Berapa \emph{ratusan} yang dimuat bilangan itu seluruhnya
(hati-hati: bukan hanya angka ratusannya --- ingatlah
\cref{ex:g3:numbers:exchange})?
\end{exercise}

\begin{exercise}[$\star$]\label{exo:g3:numbers:4}
Uraikan seperti pada \cref{def:g3:numbers:place}:
$5\,631$; \quad $2\,046$; \quad $8\,500$.
\end{exercise}

\begin{exercise}[$\star$]\label{exo:g3:numbers:5}
Salin dan lengkapi dengan $<$ atau $>$:
\[
456 \;?\; 465, \qquad
1\,203 \;?\; 989, \qquad
6\,090 \;?\; 6\,009, \qquad
9\,999 \;?\; 10\,000 .
\]
\end{exercise}

\begin{exercise}[$\star$]\label{exo:g3:numbers:6}
Urutkan dari yang terkecil ke yang terbesar: $530$; \ $3\,500$; \ $353$; \ $5\,033$;
\ $503$.
\end{exercise}

\begin{exercise}[$\star$]\label{exo:g3:numbers:7}
Pada garis bilangan berskala ratusan dari $0$ sampai $1\,000$,
letakkanlah (kira-kira) bilangan $150$, $480$, $820$ dan $990$.
\end{exercise}

\begin{exercise}[$\star$]\label{exo:g3:numbers:8}
\omterm{def:g3:numbers:evenodd}{Genap} atau \omterm{def:g3:numbers:evenodd}{ganjil}? \ $34$; \ $57$; \ $70$; \ $2\,481$; \ $6\,548$.
\end{exercise}

\begin{exercise}[$\star$]\label{exo:g3:numbers:9}
Bilanglah dengan langkah: dari $2\,970$, bilanglah lima langkah
$10$; dari $860$, empat langkah $100$; dari $9\,996$, empat langkah
$1$. Apa yang kamu perhatikan di akhir langkah yang terakhir?
\end{exercise}

\begin{exercise}[$\star\star$]\label{exo:g3:numbers:10}
Dengan memakai masing-masing angka $3$, $7$, $0$, $5$ tepat sekali,
tulislah bilangan empat angka yang terbesar, lalu yang terkecil
(sebuah bilangan tidak boleh diawali $0$).
\end{exercise}

\begin{exercise}[$\star\star$]\label{exo:g3:numbers:11}
Aku bilangan tiga angka. Angkaku $2$, $5$ dan $8$, aku \omterm{def:g3:numbers:evenodd}{genap}, dan
aku lebih besar daripada $800$. Siapakah aku?

\end{exercise}
